\documentclass[11pt,a4paper]{article}

% -------------------------------------------------
% Paquetes
% -------------------------------------------------
\usepackage[utf8]{inputenc}
\usepackage[T1]{fontenc}
\usepackage[spanish]{babel}
\usepackage{geometry}
\usepackage{enumitem}
\usepackage{hyperref}
\usepackage{xcolor}
\usepackage{listings}
\usepackage{booktabs}
\usepackage{array}
\usepackage{tabularx}
\usepackage{parskip}

\geometry{
  top=2.5cm,
  bottom=2.5cm,
  left=2.7cm,
  right=2.7cm
}

\hypersetup{
  colorlinks=true,
  linkcolor=black,
  urlcolor=blue
}

\setlist[itemize]{
  noitemsep,
  topsep=4pt
}

\setlist[enumerate]{
  noitemsep,
  topsep=4pt
}

\lstset{
  basicstyle=\ttfamily\small,
  frame=single,
  breaklines=true,
  columns=fullflexible,
  showstringspaces=false
}

\lstset{
  literate=
  {á}{{\'a}}1 {é}{{\'e}}1 {í}{{\'i}}1 {ó}{{\'o}}1 {ú}{{\'u}}1
  {Á}{{\'A}}1 {É}{{\'E}}1 {Í}{{\'I}}1 {Ó}{{\'O}}1 {Ú}{{\'U}}1
  {ñ}{{\~n}}1 {Ñ}{{\~N}}1 {¿}{{?`}}1 {¡}{{!`}}1
}

% -------------------------------------------------
% Documento
% -------------------------------------------------

\title{
  \textbf{Guía de integración Google Workspace para TI2}\\
  \large Taller de Integración II y Taller de Integración IV (TI4-16, Parte 1)
}

\author{}
\date{Segundo semestre 2026}

\begin{document}

\maketitle

% =================================================
\section{Objetivo y alcance del Sprint 1}
% =================================================

Esta guía prepara a TI2 para iniciar en Semana 2 la autenticación
institucional con Google Workspace. Es una tarea de investigación y
documentación (TI4-16). TI4 no implementa la integración en Backend, Web ni
Database en este sprint.

La Parte 1 cubre solo lo necesario para comenzar:

\begin{itemize}
  \item delimitar qué es autenticación institucional en este proyecto;
  \item explicar Google Workspace con OAuth 2.0 y OpenID Connect;
  \item explicar cómo encajan Better Auth y Convex;
  \item documentar los dominios \texttt{@alu.uct.cl} y \texttt{@uct.cl};
  \item distinguir autenticación de autorización;
  \item identificar variables, callbacks y redirects sin valores reales;
  \item dejar referencias oficiales y un checklist corto.
\end{itemize}

Queda expresamente fuera de esta tarea y del Sprint 1 la integración con
Google Calendar. Calendar se tratará más adelante como asistencia de TI4 a
TI2, con guía técnica y apoyo de implementación cuando corresponda al alcance
de agenda.

Esta guía complementa la guía de tecnologías, arquitectura y despliegue, y se
coordina con TI2-3 (proveedor y sesión), TI2-14 (logout, errores y seguridad)
y TI2-15 (pruebas y documentación).

% =================================================
\section{Delimitación: qué sí es Sprint 1}
% =================================================

Autenticación institucional en Sprint 1 significa: un usuario con cuenta
Google de la universidad inicia sesión con OAuth, el backend reconoce su
identidad y el cliente mantiene una sesión válida.

\begin{lstlisting}
Cuenta Google UCT
      |
      v
OAuth 2.0 + OpenID Connect
      |
      v
Better Auth sobre Convex entrega identidad
      |
      v
Convex valida JWT en cada función
      |
      v
Web (React + Vite + TanStack Router) mantiene sesión
\end{lstlisting}

No es Sprint 1: roles finos por recurso, organizaciones, login con RUT o
contraseña propia, doble factor, auditoría, ni Calendar. Esas piezas
pertenecen a TI2-3, TI2-14 y TI2-15, o a sprints posteriores.

El stack de referencia ya está fijado: React con Vite y TanStack Router en
Web, Convex como único backend y Better Auth como sistema de autenticación.
No se usa TanStack Start porque Convex ya es el servidor; un segundo servidor
Node solo agregaría complejidad innecesaria.

% =================================================
\section{Google Workspace, OAuth 2.0 y OpenID Connect}
% =================================================

Google Workspace es el proveedor de identidad. La universidad administra las
cuentas, Google verifica la contraseña y el segundo factor, y nuestro sistema
solo recibe una identidad verificada. Nunca vemos ni guardamos la contraseña
institucional.

El flujo estándar es Authorization Code con OpenID Connect:

\begin{enumerate}
  \item la Web redirige a \texttt{accounts.google.com} con cliente, scopes y
    \texttt{hd} cuando corresponda;
  \item el usuario autentica con Google y aprueba los scopes;
  \item Google devuelve un código a nuestro callback;
  \item Better Auth intercambia el código por tokens y obtiene el perfil
    OpenID (\texttt{openid}, \texttt{email}, \texttt{profile});
  \item se crea la sesión y Convex emite su propio JWT para las funciones.
\end{enumerate}

Los scopes del Sprint 1 son solo \texttt{openid}, \texttt{email} y
\texttt{profile}. Pedir scopes de Calendar, Drive o correo queda prohibido en
esta etapa.

En Google Cloud Console el proyecto necesita un OAuth Client de tipo Web
application con:

\begin{itemize}
  \item JavaScript origins autorizados: el origen de la Web y el Site URL de
    Convex;
  \item redirect URI autorizada: la ruta de callback de Better Auth en el Site
    URL de Convex, equivalente a \texttt{/api/auth/callback/google};
  \item pantalla de consentimiento con logo, contacto y scopes mínimos;
  \item estado de publicación coherente con el uso (interno o externo según
    defina la universidad).
\end{itemize}

Esta guía no incluye valores reales de clientes, secretos ni URLs de
deployment. Solo los nombres de las variables y la forma de las rutas.

% =================================================
\section{Consola de Google Cloud paso a paso}
% =================================================

Todo lo anterior se configura en \texttt{console.cloud.google.com}. Es la
parte más propensa a errores porque un dato mal copiado bloquea el login con
mensajes crípticos. Sigue este orden y verifica cada paso antes de pasar al
siguiente.

\subsection{1. Proyecto}

Crea o selecciona un único proyecto del equipo. Un solo proyecto es más
simple de administrar y basta para este alcance. Todo lo que sigue vive
dentro de ese proyecto: pantalla de consentimiento, credenciales y usuarios
de prueba. No es necesario habilitar APIs adicionales en la biblioteca,
porque el Sprint 1 solo usa autenticación (\texttt{openid}, \texttt{email} y
\texttt{profile}) y no consume ninguna API de Google.

\subsection{2. Pantalla de consentimiento OAuth}

Ruta: APIs y servicios, Pantalla de consentimiento. Define lo que ve el
usuario cuando Google le pide permiso.

\begin{itemize}
  \item Tipo de usuario \textbf{interno}: solo cuentas de la organización
    Workspace. No necesita usuarios de prueba y es la opción esperable si la
    universidad administra el Workspace.
  \item Tipo de usuario \textbf{externo}: cualquier cuenta Google. Mientras la
    app esté en modo de prueba, solo entran los usuarios registrados como
    test users; publicar exige verificación de Google.
  \item Datos obligatorios: nombre de la aplicación, correo de asistencia y
    dominios autorizados.
  \item Scopes: solo \texttt{openid}, \texttt{email} y \texttt{profile}. Son
    scopes no sensibles, por lo que no disparan el proceso de verificación
    extendida que exigen Calendar o Drive.
\end{itemize}

Si la universidad administra el Workspace, coordina con su TI qué tipo de
usuario corresponde y si el administrador debe permitir la app: un admin
puede bloquear apps de terceros y Google responde con
\texttt{admin\_policy\_enforced}.

\subsection{3. ID de cliente OAuth}

Ruta: APIs y servicios, Credenciales, Crear credenciales, ID de cliente
OAuth. Tipo: \textbf{Aplicación web}.

\begin{itemize}
  \item Orígenes de JavaScript autorizados: el origen de la Web en desarrollo
    (localhost con su puerto) y el origen de producción en Vercel.
  \item URIs de redirección autorizadas: la ruta de callback de Better Auth
    bajo el Site URL de Convex, con forma equivalente a
    \texttt{https://<deployment>.convex.site/api/auth/callback/google}.
    Debe coincidir exacta: esquema, mayúsculas y slash final. Cualquier
    diferencia devuelve \texttt{redirect\_uri\_mismatch}.
  \item Al crear el cliente, copia el Client ID y el Client Secret. El secreto
    se muestra una sola vez. Guárdalos fuera del repositorio y llévalos al
    deployment solo con \texttt{convex env set}.
\end{itemize}

\subsection{4. Errores típicos y su causa}

\begin{description}
  \item[\texttt{redirect\_uri\_mismatch}:] el callback no está registrado o
    difiere en un carácter. Revisa la lista de URIs en Credenciales.
  \item[\texttt{org\_internal}:] cliente de tipo interno con una cuenta fuera
    de la organización. Entra con cuenta institucional o revisa el tipo de
    usuario.
  \item[\texttt{admin\_policy\_enforced}:] el admin del Workspace bloquea apps
    de terceros. Lo resuelve el administrador, no el código.
  \item[\texttt{invalid\_client}:] secreto incorrecto o cliente eliminado.
    Regenera el secreto y actualiza el deployment.
  \item[Acceso bloqueado en modo de prueba:] la cuenta no es test user o la
    app externa aún no se publica. Agrega la cuenta o avanza la publicación.
\end{description}

% =================================================
\section{Dominios institucionales}
% =================================================

El proyecto distingue dos dominios:

\begin{description}
  \item[\texttt{@alu.uct.cl}:] estudiantes.
  \item[\texttt{@uct.cl}:] docentes, funcionarios y administración.
\end{description}

Better Auth soporta la opción \texttt{hd} del proveedor Google para exigir un
hosted domain de Workspace. El parámetro \texttt{hd} sugiere el dominio en la
pantalla de Google y, configurado en el proveedor, hace que Better Auth
rechace tokens cuyo claim \texttt{hd} no coincida.

El patrón probado en cuestionarios-uct usa dos inicios lógicos sobre el mismo
OAuth Client de Google, uno con \texttt{hd} de estudiantes y otro con
\texttt{hd} de docentes, más una validación dura del sufijo del email en el
mapeo de perfil para el rol docente. El parámetro \texttt{hd} por sí solo es
una sugerencia de UX; la garantía real es validar el dominio en el servidor.

\begin{lstlisting}
socialProviders: {
  google: {
    clientId: process.env.GOOGLE_CLIENT_ID,
    clientSecret: process.env.GOOGLE_CLIENT_SECRET,
    hd: "alu.uct.cl",
  },
}
\end{lstlisting}

Para dos poblaciones, TI2 puede optar por un proveedor con \texttt{hd} único
si el sprint decide un solo dominio inicial, o por la variante de dos
proveedores lógicos ya validada en cuestionarios-uct. Lo importante es que la
decisión quede en el servidor (Better Auth sobre Convex), nunca solo en el
frontend.

% =================================================
\section{Cómo encajan Better Auth y Convex}
% =================================================

Better Auth administra registro, login social, sesión y logout. Convex aporta
ejecución, base de datos y verificación del JWT en cada función. La
integración oficial (\texttt{@convex-dev/better-auth}) hace que Better Auth
corra como rutas HTTP dentro del deployment Convex y persista sus tablas en
el componente \texttt{betterAuth}.

\begin{lstlisting}
Web (Vite SPA)
  |
  +-- authClient.signIn.social({ provider: "google" })
  |
  v
https://<deployment>.convex.site/api/auth/*
  (Better Auth sobre Convex)
  |
  +-- Google OAuth + callback
  +-- tablas user, session, account, verification
  |
  v
Convex functions verifican ctx.auth.getUserIdentity()
\end{lstlisting}

Los archivos del backend son siempre los mismos:

\begin{lstlisting}
convex/
|
+-- auth.config.ts       proveedor JWT de Better Auth
+-- convex.config.ts     registra el componente betterAuth
+-- http.ts              expone /api/auth/* en Convex Site
+-- auth.ts              createAuth: baseURL, adapter,
|                       socialProviders.google, trustedOrigins
+-- schema.ts            solo dominio; auth vive en el componente
\end{lstlisting}

Fragmento de forma (sin valores):

\begin{lstlisting}
auth.config.ts:
  providers: [getAuthConfigProvider()]

convex.config.ts:
  app.use(betterAuth)

http.ts:
  authComponent.registerRoutes(http, createAuth)

auth.ts:
  baseURL: process.env.BETTER_AUTH_URL
  database: authComponent.adapter(ctx)
  socialProviders: { google: { clientId, clientSecret } }
  plugins: [convex(), crossDomain()]
\end{lstlisting}

En la Web Vite (SPA, sin Start) el cliente es:

\begin{lstlisting}
import { createAuthClient } from "better-auth/react";
import { convexClient, crossDomainClient } from
  "@convex-dev/better-auth/client/plugins";

export const authClient = createAuthClient({
  baseURL: import.meta.env.VITE_CONVEX_SITE_URL,
  plugins: [convexClient(), crossDomainClient()],
});
\end{lstlisting}

Y el provider en la raíz:

\begin{lstlisting}
const convex = new ConvexReactClient(
  import.meta.env.VITE_CONVEX_URL
);

<ConvexBetterAuthProvider client={convex} authClient={authClient}>
  <App />
</ConvexBetterAuthProvider>
\end{lstlisting}

El login social desde TanStack Router es una llamada al cliente, con
\texttt{callbackURL} y \texttt{errorCallbackURL} del frontend:

\begin{lstlisting}
await authClient.signIn.social({
  provider: "google",
  callbackURL: "/",
  errorCallbackURL: "/login",
});
\end{lstlisting}

La sesión reactiva se lee con \texttt{authClient.useSession()}. La verdad
autoritativa, sin embargo, siempre está en el backend con
\texttt{ctx.auth.getUserIdentity()} o
\texttt{authComponent.getAuthUser(ctx)}. El patrón verificado en ponderapp y
reportevecino es: nunca aceptar un \texttt{userId} desde los argumentos para
autorizar; la identidad viene del JWT.

% =================================================
\section{Autenticación no es autorización}
% =================================================

\begin{description}
  \item[Autenticación:] quién es el usuario. La resuelve Better Auth con
    Google y la entrega como identidad.
  \item[Autorización:] qué puede hacer. La resuelven los casos de uso y el
    dominio, en el backend, a partir de esa identidad.
\end{description}

\begin{lstlisting}
Infraestructura
    |
    +--> Better Auth entrega identidad (quién)
    |
    v
Aplicación y Dominio evalúan permiso (qué puede)
    |
    v
Convex ejecuta o rechaza
\end{lstlisting}

Ocultar un botón en React no es autorización. Toda regla (rol, recurso,
pertenencia) debe verificarse en la función Convex antes de leer o mutar.
Los roles finos y la matriz de permisos se definen en TI2-3 y se endurecen en
TI2-14; en Sprint 1 basta con distinguir usuario autenticado de anónimo y
dejar el punto de extensión para roles.

% =================================================
\section{Variables, callbacks y redirects}
% =================================================

Solo nombres. Ningún valor real pertenece a esta guía ni al repositorio.

\begin{table}[h]
  \centering
  \begin{tabularx}{\textwidth}{>{\bfseries}l X}
    \toprule
    Variable & Uso \\
    \midrule
    \texttt{GOOGLE\_CLIENT\_ID} & OAuth Client de Google, solo servidor \\
    \texttt{GOOGLE\_CLIENT\_SECRET} & Secreto OAuth, solo servidor \\
    \texttt{BETTER\_AUTH\_SECRET} & Firma de sesiones, solo servidor \\
    \texttt{BETTER\_AUTH\_URL} & Base URL de Better Auth en Convex \\
    \texttt{BETTER\_AUTH\_TRUSTED\_ORIGINS} & Orígenes extra permitidos \\
    \texttt{VITE\_CONVEX\_URL} & Cliente Convex, pública \\
    \texttt{VITE\_CONVEX\_SITE\_URL} & Base de \texttt{/api/auth/*}, pública \\
    \bottomrule
  \end{tabularx}
\end{table}

Reglas:

\begin{itemize}
  \item las secretas se configuran con \texttt{convex env set} en el
    deployment, nunca con prefijo \texttt{VITE\_} ni en el bundle;
  \item las públicas \texttt{VITE\_} son visibles para cualquiera; solo
    localizan servicios;
  \item \texttt{trustedOrigins} debe incluir el origen de desarrollo Vite
    (puerto de \texttt{dev}), el origen de producción en Vercel y el Site URL
    de Convex; se extiende por entorno con
    \texttt{BETTER\_AUTH\_TRUSTED\_ORIGINS};
  \item el callback Google es la ruta \texttt{/api/auth/callback/google} bajo
    el Convex Site URL; el frontend solo define \texttt{callbackURL} y
    \texttt{errorCallbackURL} de navegación post-login.
\end{itemize}

Archivos de ejemplo sin secretos: \texttt{apps/web/.env.example} con
\texttt{VITE\_CONVEX\_URL} y \texttt{VITE\_CONVEX\_SITE\_URL}; locales en
\texttt{.env.local} fuera de Git.

% =================================================
\section{Checklist para que TI2 comience en Semana 2}
% =================================================

\begin{enumerate}
  \item Completar la sección de Consola de Google Cloud: proyecto, pantalla de
    consentimiento (interno o externo), OAuth Client Web con redirect al Site
    URL de Convex y scopes \texttt{openid}, \texttt{email}, \texttt{profile}.
  \item Fijar decisión de dominios: un \texttt{hd} inicial o los dos
    (\texttt{alu.uct.cl} y \texttt{uct.cl}) con validación de sufijo en el
    servidor.
  \item Levantar \texttt{convex/auth.config.ts}, \texttt{convex.config.ts},
    \texttt{http.ts} y \texttt{auth.ts} según la estructura oficial.
  \item Configurar secretos solo con \texttt{convex env set} y publicar
    \texttt{.env.example} sin valores.
  \item Conectar la SPA Vite con \texttt{authClient}, \texttt{convexClient},
    \texttt{crossDomainClient} y \texttt{ConvexBetterAuthProvider}.
  \item Implementar login y logout con \texttt{signIn.social} y
    \texttt{signOut}, más lectura con \texttt{useSession}.
  \item Verificar en una query que \texttt{ctx.auth.getUserIdentity()} distingue
    autenticado de anónimo.
  \item Registrar decisiones y dudas como comentarios en TI2-3 para coordinar
    con TI4.
\end{enumerate}

% =================================================
\section{Referencias oficiales}
% =================================================

\begin{itemize}
  \item Better Auth, Google social provider y opción \texttt{hd}:
    \url{https://better-auth.com/docs/authentication/google}
  \item Better Auth, uso de \texttt{signIn.social}:
    \url{https://better-auth.com/docs/basic-usage}
  \item Better Auth con Convex (guía marco React):
    \url{https://labs.convex.dev/better-auth/framework-guides/react}
  \item Better Auth con Convex (integración):
    \url{https://better-auth.com/docs/integrations/convex}
  \item Convex, cliente React:
    \url{https://docs.convex.dev/client/react/overview}
  \item Convex en hosting Vercel:
    \url{https://docs.convex.dev/production/hosting/vercel}
  \item Google Identity, OAuth 2.0 y OpenID Connect:
    \url{https://developers.google.com/identity/openid-connect/openid-connect}
  \item Google Cloud, configuración de pantalla de consentimiento y Client ID:
    \url{https://developers.google.com/identity/protocols/oauth2}
  \item Google Cloud, página de clientes OAuth (crear credenciales):
    \url{https://console.developers.google.com/auth/clients}
  \item Google Cloud, tipos de usuario y pantalla de consentimiento:
    \url{https://support.google.com/cloud/answer/10311615}
  \item TanStack Router con Vite:
    \url{https://tanstack.com/router/latest/docs/installation/with-vite}
\end{itemize}

Patrones internos verificados: cuestionarios-uct
(\texttt{apps/web/src/server/auth/index.ts},
  \texttt{src/lib/auth/client.ts},
\texttt{src/lib/constants.ts}) para Google con \texttt{hd} y validación por
dominio; ponderapp (\texttt{apps/backend/convex/auth.ts},
\texttt{apps/web/src/lib/auth-client.ts}) y reportevecino
(\texttt{apps/backend/convex/auth.ts},
  \texttt{apps/backend/convex/access.ts},
\texttt{apps/dashboard/src/integrations/convex/provider.tsx}) para Better Auth
sobre Convex y verificación de sesión en backend.

\end{document}
